% ch3_d (书页 131-142 / p146-p157)
%--C-TAIL--
%JOIN% 为
% ---- 书页 131 / p146 ----
\begin{equation*}
G_{\mu\nu} = R_{\mu\nu} - \frac{1}{2} R g_{\mu\nu}.
\tag{3.151}
\end{equation*}

在四维时，可以把 Einstein 张量看作 Ricci 张量的一个无迹反转（trace-reversed）版本。于是我们看到，两次缩并的 Bianchi 恒等式 (3.150) 等价于
\begin{equation*}
\nabla^\mu G_{\mu\nu} = 0.
\tag{3.152}
\end{equation*}

Einstein 张量由于 Ricci 张量和度规的对称性而是对称的，它在广义相对论中将具有极大的重要性。

此时我们应当停顿一下，把我们发展出来的这套形式体系同我们对曲率的直观概念作一番对比。遗憾的是，我们的直觉已被一个事实所污染：我们习惯于思考嵌在我们生活的（近乎）欧几里得空间中的一维和二维空间。例如，我们认为直线没有曲率，而圆（$S^1$）是弯曲的。然而，根据式 (3.137)，在一、二、三、四维中，Riemann 张量分别有 0、1、6 和 20 个独立分量。（在这些例子中我们关于曲率所说的一切，指的都是与 Christoffel 联络从而与度规相联系的曲率。）因此，像 $S^1$ 这样的一个一维空间，按我们的定义不可能具有任何曲率。这个表面上的矛盾源于如下事实：我们对曲率的直观概念依赖于流形的外在几何（extrinsic geometry），它刻画的是一个空间如何嵌入某个更大的空间；而 Riemann 曲率则是一个空间的内蕴几何（intrinsic geometry）的属性，它可以被局限于流形之上的观者测量出来。生活在圆上、无法接触更大世界的生命，必然认为他们生活在一个平直的几何中——例如，在那里根本不存在一个非退化的无穷小回路，可以让我们沿它平行移动一个矢量并使其回来时相对原来的方位有所转动。在附录 D 中讨论的外曲率，当我们想描述时空的子流形时偶尔有用；但大多数时候我们感兴趣的是时空自身的内蕴几何，它不依赖任何嵌入。

我们可以用二维的一个例子进一步说明内蕴/外在的差别，在二维中曲率只有一个独立分量。事实上，关于曲率的全部信息都包含在 Ricci 标量曲率这单个分量之中。考虑一个环面（torus），如图 3.7 所示，它可以看作平面上一个把对边两两等同起来的正方形区域（拓扑上就是 $S^1\times S^1$）。虽然嵌在三维中的环面从我们的观点看起来是弯曲的，但很清楚，我们可以给环面放置一个度规，使其分量在某个适当的坐标系中为常数——只需把它摊开铺平，使用平面的欧几里得度规 $\mathrm{d}s^2 = \mathrm{d}x^2 + \mathrm{d}y^2$ 即可。在这个度规下，环面是平直的。
% ---- 书页 132 / p147 ----
%FIG{3.7}: {把环面看作平直空间中一个对边被等同起来的正方形。}
同样，没有什么能阻止我们引入另一个度规使环面不平直；但我们要强调的要点是，它在某个度规之下可以被弄平。每当我们把一个流形嵌入更大的空间，该流形都会从它所嵌入的背景那里继承一个“诱导度规”（induced metric），如附录 A 中所讨论的。我们在这里的要点是：嵌在平直三维欧几里得空间中的环面将具有弯曲的诱导度规，但我们仍然可以选择给它放置另一个不同的度规，使内蕴几何是平直的。

让我们转向一个曲率并不为零的简单例子。我们已经谈论过二维球面 $S^2$，其度规为
\begin{equation*}
ds^2 = a^2(\mathrm{d}\theta^2 + \sin^2\theta\,\mathrm{d}\phi^2),
\tag{3.153}
\end{equation*}
其中 $a$ 是球的半径。如果我们的球嵌在 $\mathbf{R}^3$ 中，它实际上就是半径；但即使没有任何嵌入，我们仍然可以称它为半径。生活在球面上的二维居民想要算出 $a$，可以测量球的面积，除以 $4\pi$，再取平方根；用“半径”一词来指称这个量仅仅是为了方便。我们还应当指出，球的概念有时在更弱的拓扑意义下使用，不假定任何具体的度规；我们使用的这个度规称为\emph{圆度规}（round metric）。不必细究细节，式 (3.153) 的非零联络系数为
\begin{align*}
\Gamma^\theta_{\phi\phi} &= -\sin\theta\cos\theta\\
\Gamma^\phi_{\theta\phi} = \Gamma^\phi_{\phi\theta} &= \cot\theta.
\tag{3.154}
\end{align*}
我们来计算一个颇有希望的 Riemann 张量分量：
\begin{align*}
R^\theta{}_{\phi\theta\phi} &= \partial_\theta\Gamma^\theta_{\phi\phi} - \partial_\phi\Gamma^\theta_{\theta\phi} + \Gamma^\theta_{\theta\lambda}\Gamma^\lambda_{\phi\phi} - \Gamma^\theta_{\phi\lambda}\Gamma^\lambda_{\theta\phi}\\
&= (\sin^2\theta - \cos^2\theta) - (0) + (0) - (-\sin\theta\cos\theta)(\cot\theta)\\
&= \sin^2\theta.
\tag{3.155}
\end{align*}
% ---- 书页 133 / p148 ----
这个记号显然不完美，因为希腊字母 $\lambda$ 是一个被求和的哑指标，而希腊字母 $\theta$ 和 $\phi$ 则代表具体的坐标。下降一个指标，我们得到
\begin{align*}
R_{\theta\phi\theta\phi} &= g_{\theta\lambda}R^\lambda{}_{\phi\theta\phi}\\
&= g_{\theta\theta}R^\theta{}_{\phi\theta\phi}\\
&= a^2\sin^2\theta.
\tag{3.156}
\end{align*}
容易验证，Riemann 张量的所有分量要么为零，要么通过对称性与这个分量相联系。我们可以继续计算 Ricci 张量，利用 $R_{\mu\nu} = g^{\alpha\beta}R_{\alpha\mu\beta\nu}$。我们得到
\begin{align*}
R_{\theta\theta} &= g^{\phi\phi}R_{\phi\theta\phi\theta} = 1\\
R_{\theta\phi} &= R_{\phi\theta} = 0\\
R_{\phi\phi} &= g^{\theta\theta}R_{\theta\phi\theta\phi} = \sin^2\theta.
\tag{3.157}
\end{align*}
Ricci 标量曲率同样直截了当：
\begin{equation*}
R = g^{\theta\theta}R_{\theta\theta} + g^{\phi\phi}R_{\phi\phi} = \frac{2}{a^2}.
\tag{3.158}
\end{equation*}

因此，对二维流形完全刻画了曲率的 Ricci 标量曲率，在整个二维球面上是一个常数。如果我们扰动了这个度规（物理上对应球面上的凹凸不平），情形就不再如此。注意，标量曲率随球半径的增大而减小。即使在更一般的语境中，我们有时也会提到流形的“曲率半径”（radius of curvature），用它来提供一个曲率变化的长度尺度；曲率半径越大，曲率本身越小。

\section{对称性与 Killing 矢量}

真实世界是个乱糟糟的地方，我们别指望找到一个度规，能以完美的精度描述我们实际的宇宙，乃至其任何一个小的部分。我们只能转而借助各种适合于所研究物理情形的近似来为时空建模。例如，恒星或行星外部的几何，可以近似到某种精度地看作球对称的，即使真实情形包含对这种对称性的小偏离——这些偏离可以在以后作为微扰加进来。

广义相对论在这一点上同其他物理学领域并无不同，也对具有对称性的解特别感兴趣。事实上，这类性质在广义相对论中甚至可能比在比如电磁学中更加关键，因为 Einstein 方程（下一章讨论）的非线性本质使得寻找任何精确解都很困难。然而，在弯曲时空的语境下，对于“对称性”究竟意味着什么，我们需要比平时更加小心。本节将发展一些研究对称性的有用工具；更深入的探讨可以在附录 B 中找到。
% ---- 书页 134 / p149 ----
如果几何在一个把 $M$ 映到其自身的变换下不变，我们就认为流形 $M$ 具有一种对称性；也就是说，度规在某种意义上从一点到另一点都是相同的。事实上，不同的张量场可以拥有不同的对称性；度规的对称性称为\textbf{等距变换}（isometries）。有时等距变换的存在是显而易见的；例如考虑四维 Minkowski 空间，
\begin{equation*}
ds^2 = \eta_{\mu\nu}\mathrm{d}x^\mu\mathrm{d}x^\nu = -\mathrm{d}t^2 + \mathrm{d}x^2 + \mathrm{d}y^2 + \mathrm{d}z^2.
\tag{3.159}
\end{equation*}
我们知道这个空间有若干等距变换；它们包括平移（$x^\mu \to x^\mu + a^\mu$，其中 $a^\mu$ 固定）和洛伦兹变换（$x^\mu \to \Lambda^\mu{}_\nu x^\nu$，其中 $\Lambda^\mu{}_\nu$ 是一个洛伦兹变换矩阵）。度规在平移下不变这一点，由一个简单的事实即可立即看出：度规系数 $\eta_{\mu\nu}$ 不依赖于各个坐标函数 $x^\mu$。事实上，只要对某个固定的 $\sigma_*$（但对所有的 $\mu$ 和 $\nu$）都有 $\partial_{\sigma_*}g_{\mu\nu}=0$，就会存在沿 $x^{\sigma_*}$ 方向平移的对称性：
\begin{equation*}
\partial_{\sigma_*}g_{\mu\nu} = 0 \;\;\Rightarrow\;\; x^{\sigma_*} \to x^{\sigma_*} + a^{\sigma_*}\ \text{是一个对称性.}
\tag{3.160}
\end{equation*}

细心的读者会注意到，我们还没有精确定义我们所说的对称性是什么；粗略地说，我们设想度规在某个变换下不变，但精确的含义要到附录 B 中才展开。另外，式 (3.160) 中的蕴含箭头只是单向的，若能有一个干净的判据来判定一个给定的变换何时算作对称性，那就更好了；这一点很快就会讲到。

式 (3.160) 这种形式的等距变换，对测地线方程所描述的检验粒子的运动有直接的推论。回忆式 (3.61)，测地线方程可以用四维动量 $p^\mu = mU^\mu$（至少对类时路径有效）写成
\begin{equation*}
p^\lambda\nabla_\lambda p^\mu = 0.
\tag{3.161}
\end{equation*}
由度规相容性，我们可以自由地下降指标 $\mu$，然后把协变导数展开，得到
\begin{equation*}
p^\lambda\partial_\lambda p_\mu - \Gamma^\sigma_{\lambda\mu}p^\lambda p_\sigma = 0.
\tag{3.162}
\end{equation*}
第一项告诉我们动量的分量沿路径如何变化，
\begin{equation*}
p^\lambda\partial_\lambda p_\mu = m\frac{dx^\lambda}{d\tau}\partial_\lambda p_\mu = m\frac{dp_\mu}{d\tau},
\tag{3.163}
\end{equation*}
% ---- 书页 135 / p150 ----
而第二项则是
\begin{align*}
\Gamma^\sigma_{\lambda\mu}p^\lambda p_\sigma &= \frac{1}{2}g^{\sigma\nu}(\partial_\lambda g_{\mu\nu} + \partial_\mu g_{\nu\lambda} - \partial_\nu g_{\lambda\mu})p^\lambda p_\sigma \tag{3.164}\\
&= \frac{1}{2}(\partial_\lambda g_{\mu\nu} + \partial_\mu g_{\nu\lambda} - \partial_\nu g_{\lambda\mu})p^\lambda p^\nu \tag{3.165}\\
&= \frac{1}{2}(\partial_\mu g_{\nu\lambda})p^\lambda p^\nu,
\tag{3.166}
\end{align*}
其中我们利用了 $p^\lambda p^\nu$ 的对称性从第二行过渡到第三行。于是，在尚未对对称性作任何假设的情况下，我们就看到测地线方程可以写成
\begin{equation*}
m\frac{dp^\mu}{d\tau} = \frac{1}{2}(\partial_\mu g_{\nu\lambda})p^\lambda p^\nu.
\tag{3.167}
\end{equation*}
因此，如果所有的度规系数都不依赖于坐标 $x^{\sigma_*}$，我们就发现这个等距变换蕴含动量分量 $p_{\sigma_*}$ 是该运动的一个守恒量：
\begin{equation*}
\partial_{\sigma_*}g_{\mu\nu} = 0 \;\;\Rightarrow\;\; \frac{dp_{\sigma_*}}{d\tau} = 0.
\tag{3.168}
\end{equation*}

这个结果沿任何测地线都成立，尽管我们只是对类时测地线导出了它。等距变换所蕴含的守恒量，对研究弯曲背景中检验粒子的运动极为有用。

当然，虽然度规分量对一个或多个坐标的独立性蕴含等距变换的存在，反过来却未必成立。例如，洛伦兹变换下的对称性并不表现为 $\eta_{\mu\nu}$ 对任何坐标的独立性；事实上，在四维中共有四类平移和六类洛伦兹变换，总计十个，这显然大于度规可能与之无关的坐标数目。更何况，我们完全可以变换到一个复杂的坐标系，在其中连平移对称性都看不出来。这样的坐标变换会改变度规的分量，却不改变底层的几何——而对称性真正刻画的正是后者。显然，我们需要一个更系统的程序。

我们可以发展出这样一个程序：把式 (3.168) 右边的方程——它表达动量的某个分量为常数——改写成一种协变性更加明显的形式。如果 $x^{\sigma_*}$ 就是 $g_{\mu\nu}$ 与之无关的那个坐标，让我们考虑矢量 $\partial_{\sigma_*}$，把它记作 $K$：
\begin{equation*}
K = \partial_{\sigma_*},
\tag{3.169}
\end{equation*}
用分量记号来写，它等价于
\begin{equation*}
K^\mu = (\partial_{\sigma_*})^\mu = \delta^\mu_{\sigma_*}.
\tag{3.170}
\end{equation*}
我们说矢量 $K^\mu$ 生成（generate）这个等距变换；这意味着，几何在其下不变的变换无穷小地表达为沿 $K^\mu$ 方向的一个移动。
% ---- 书页 136 / p151 ----
同样，这个概念在附录 B 中有更充分的展开。用这个矢量来表示，看上去不具协变形式的量 $p_{\sigma_*}$ 就简单地是
\begin{equation*}
p_{\sigma_*} = K^\nu p_\nu = K_\nu p^\nu.
\tag{3.171}
\end{equation*}
另一方面，这个（标量）量沿路径的恒定性，等价于说它沿测地线的方向导数为零：
\begin{equation*}
\frac{dp_{\sigma_*}}{d\tau} = 0 \;\;\Longleftrightarrow\;\; p^\mu\nabla_\mu(K_\nu p^\nu) = 0.
\tag{3.172}
\end{equation*}
把右边的表达式展开，我们得到
\begin{align*}
p^\mu\nabla_\mu(K_\nu p^\nu) &= p^\mu K_\nu\nabla_\mu p^\nu + p^\mu p^\nu\nabla_\mu K_\nu\\
&= p^\mu p^\nu\nabla_\mu K_\nu\\
&= p^\mu p^\nu\nabla_{(\mu}K_{\nu)},
\tag{3.173}
\end{align*}
其中在第二行我们引用了测地线方程（$p^\mu\nabla_\mu p^\nu = 0$）。在第三行我们利用了 $p^\mu p^\nu$ 在 $\mu$ 和 $\nu$ 上自动对称这一事实，因此只有 $\nabla_\mu K_\nu$ 的对称部分才可能给出贡献。于是我们得出结论：任何满足 $\nabla_{(\mu}K_{\nu)} = 0$ 的矢量 $K_\mu$ 都蕴含 $K_\nu p^\nu$ 沿测地线轨道守恒：
\begin{equation*}
\nabla_{(\mu}K_{\nu)} = 0 \;\;\Rightarrow\;\; p^\mu\nabla_\mu(K_\nu p^\nu) = 0.
\tag{3.174}
\end{equation*}

左边的方程称为\textbf{Killing 方程}（Killing's equation），满足它的矢量场称为\textbf{Killing 矢量场}（Killing vector fields，或简称 Killing 矢量）。你可以自己验证：如果度规不依赖于某个坐标 $x^{\sigma_*}$，矢量 $\partial_{\sigma_*}$ 将满足 Killing 方程。事实上，如果一个矢量 $K^\mu$ 满足 Killing 方程，那么总能找到一个坐标系使得 $K = \partial_{\sigma_*}$；但一般说来，我们无法找到一个坐标系使所有 Killing 矢量都同时取这种形式，而且矢量要满足 Killing 方程也并不需要这种形式。

正如我们将在附录 B 中研究的，流形上的 Killing 矢量场与该流形上度规的连续对称性一一对应。每个 Killing 矢量都蕴含着与测地线运动相联系的守恒量的存在。这一点可以从物理上理解：按定义，度规沿 Killing 矢量的方向是不变的。粗略地说，自由粒子在这个方向上不会感到任何力，因而其动量在这个方向上的分量将相应地守恒。事实上，我们用来证明“若 $\nabla_{(\mu}K_{\nu)} = 0$ 则 $K_\nu p^\nu$ 沿测地线守恒”的那类逻辑，可以推广到更多的指标：\textbf{Killing 张量}（Killing tensor）是一个对称的 $\ell$ 指标张量 $K_{\nu_1\ldots\nu_\ell}$，它满足 Killing 方程的显而易见的推广，并相应地通过与 $\ell$ 份动量缩并而导致守恒量：
% ---- 书页 137 / p152 ----
\begin{equation*}
\nabla_{(\mu}K_{\nu_1\ldots\nu_\ell)} = 0 \;\;\Rightarrow\;\; p^\mu\nabla_\mu(K_{\nu_1\ldots\nu_\ell}p^{\nu_1}\cdots p^{\nu_\ell}) = 0.
\tag{3.175}
\end{equation*}

Killing 张量的简单例子有度规本身，以及 Killing 矢量的对称化张量积。Killing 张量与时空中对称性的联系并不简单，但它们将简化我们对转动黑洞和膨胀宇宙的分析。

Killing 矢量的导数可以通过下式与 Riemann 张量联系起来
\begin{equation*}
\nabla_\mu\nabla_\sigma K^\rho = R^\rho{}_{\sigma\mu\nu}K^\nu,
\tag{3.176}
\end{equation*}
正如习题中要求你们证明的那样。把这个表达式缩并，得到
\begin{equation*}
\nabla_\mu\nabla_\sigma K^\mu = R_{\sigma\nu}K^\nu.
\tag{3.177}
\end{equation*}
这些关系连同 Bianchi 恒等式与 Killing 方程，足以证明 Ricci 标量曲率沿 Killing 矢量场的方向导数为零，
\begin{equation*}
K^\lambda\nabla_\lambda R = 0.
\tag{3.178}
\end{equation*}

最后这个事实从另一个侧面反映了这样的观念：几何沿 Killing 矢量场不发生变化。

除了为单个粒子的运动导致守恒量之外，类时 Killing 矢量的存在还使我们能够为整个时空定义一个守恒的能量。给定一个 Killing 矢量 $K_\nu$ 和一个守恒的能量-动量张量 $T_{\mu\nu}$，我们可以构造一个流
\begin{equation*}
J^\mu_T = K_\nu T^{\mu\nu}
\tag{3.179}
\end{equation*}
它自动守恒，
\begin{align*}
\nabla_\mu J^\mu_T &= (\nabla_\mu K_\nu)T^{\mu\nu} + K_\nu(\nabla_\mu T^{\mu\nu})\\
&= 0.
\tag{3.180}
\end{align*}
第一项由于 Killing 方程而为零（因为上指标的对称性恰好使下指标自动对称化），第二项则由于 $T_{\mu\nu}$ 的守恒而为零。如果 $K_\nu$ 是类时的，我们可以在一个类空超曲面 $\Sigma$ 上积分来定义总能量，
\begin{equation*}
E_T = \int_\Sigma J^\mu_T n_\mu \sqrt{\gamma}\,\mathrm{d}^3x,
\tag{3.181}
\end{equation*}
其中 $\gamma_{ij}$ 是 $\Sigma$ 上的诱导度规，$n_\mu$ 是 $\Sigma$ 的法矢量。在附录 E 中我们将讨论超曲面上的积分，特别是 Stokes 定理；正如那里所解释的，$E_T$ 在任何类空超曲面上积分都得到相同的值，因而是守恒的。
% ---- 书页 138 / p153 ----
这个结果与我们在第 3.5 节中的讨论相当契合：在那里我们发现，在一个膨胀的宇宙中总能量通常不守恒；膨胀意味着度规随时间变化，因此在这个方向上没有等距变换。当存在类时 Killing 矢量时，我们可以把度规写成不依赖于类时坐标的形式，而诺特定理便蕴含一个守恒的能量。类似地，类空 Killing 矢量可以用来构造守恒的动量（或角动量）。

虽然在任意给定的时空中实际求解 Killing 方程可能简单也可能不简单，但通过直接观察写下某些 Killing 矢量往往是可能的。（当然，一个一般的度规可能根本没有任何 Killing 矢量，但为了保持简单，我们经常处理具有高度对称性的度规。）例如，在具有度规 $\mathrm{d}s^2 = \mathrm{d}x^2 + \mathrm{d}y^2 + \mathrm{d}z^2$ 的 $\mathbf{R}^3$ 中，度规分量对 $x$、$y$、$z$ 的无关性立即给出三个 Killing 矢量：
\begin{align*}
X^\mu &= (1,\, 0,\, 0)\\
Y^\mu &= (0,\, 1,\, 0)\\
Z^\mu &= (0,\, 0,\, 1).
\tag{3.182}
\end{align*}
这些显然代表三个平移。$\mathbf{R}^3$ 中还有三个转动对称性，它们就不那么简单了。为了找到它们，设想先变换到球坐标，
\begin{align*}
x &= r\sin\theta\cos\phi\\
y &= r\sin\theta\sin\phi\\
z &= r\cos\theta,
\tag{3.183}
\end{align*}
此时度规取如下形式
\begin{equation*}
ds^2 = \mathrm{d}r^2 + r^2\mathrm{d}\theta^2 + r^2\sin^2\theta\,\mathrm{d}\phi^2.
\tag{3.184}
\end{equation*}
现在这个度规（还是\emph{同一个}度规，只是换了一个坐标系）明显不依赖于 $\phi$。因此我们知道 $R = \partial_\phi$ 是一个 Killing 矢量。变换回笛卡尔坐标，它就成为
\begin{equation*}
R = -y\partial_x + x\partial_y.
\tag{3.185}
\end{equation*}
于是 $R^\mu$ 的笛卡尔分量就是 $(-y,\, x,\, 0)$。由于这代表绕 $z$ 轴的一个转动，不难猜出所有三个转动 Killing 矢量的分量：
\begin{align*}
R^\mu &= (-y,\; x,\; 0)\\
S^\mu &= (\phantom{-}z,\; 0,\; -x)\\
T^\mu &= (\;0,\; -z,\; y),
\tag{3.186}
\end{align*}
% ---- 书页 139 / p154 ----
分别代表绕 $z$、$y$、$x$ 轴的转动。你可以自己验证，这些矢量确实求解 Killing 方程。整体的正负号无关紧要，因为一个 Killing 矢量加负号仍然是 Killing 矢量。

这个练习直接引出二维球面 $S^2$ 的 Killing 矢量，其度规为
\begin{equation*}
ds^2 = \mathrm{d}\theta^2 + \sin^2\theta\,\mathrm{d}\phi^2.
\tag{3.187}
\end{equation*}
由于这个球可以看作 $\mathbf{R}^3$ 中到原点距离为 1 的点的轨迹，而各个转动的 Killing 矢量都把这样一个球转成它自身，所以它们也表示 $S^2$ 的对称性。为了得到这些矢量的明显的坐标基表示，我们先把三维矢量 (3.186) 变换到球坐标 $x^{\mu'} = (r, \theta, \phi)$。一番直截了当的计算给出
\begin{align*}
R &= \partial_\phi\\
S &= \cos\phi\,\partial_\theta - \cot\theta\sin\phi\,\partial_\phi\\
T &= -\sin\phi\,\partial_\theta - \cot\theta\cos\phi\,\partial_\phi.
\tag{3.188}
\end{align*}
注意，这里没有沿 $\partial_r$ 的分量，这对一个转动的等距变换来说正合情理。因此，$\mathbf{R}^3$ 中三个转动 Killing 矢量的表达式 (3.188)，与球坐标下 $S^2$ 的那些表达式完全相同。

在 $n\geq 2$ 维时，Killing 矢量的数目可以多于维数。这是因为一组 Killing 矢量场可以线性无关，尽管在流形上任何一点处，该点上的这些矢量是线性相关的。证明 Killing 矢量以\emph{常数}系数的线性组合仍是 Killing 矢量是很容易的（所以你应该自己动手做一下；此时该线性组合不再算作一个独立的 Killing 矢量），但对于在流形上变化的系数，这一般不再成立。你还可以证明，两个 Killing 矢量场的对易子是一个 Killing 矢量场；知道这一点非常有用，但对易子给出的矢量场可能不是线性无关的（也可能干脆为零）。因此，求出一个度规的所有 Killing 矢量这个问题颇为棘手，因为什么时候该停止寻找并不总是清楚的。

\section{最大对称空间}

一个空间究竟能有多对称？具有最高可能对称度的空间的一个例子，是带有平直欧几里得度规的 $\mathbf{R}^n$。让我们从它们在某个固定点 $p$ 的邻域内做什么这个角度，来考察这个空间的等距变换——我们知道它们就是 $n$ 维中的平移和转动。平移是移动该点的那些变换；可以有 $n$ 个独立的轴沿着移动它，所以平移总共 $n$ 个。以 $p$ 为中心的转动是使 $p$ 保持不变的那些变换；可以把它们看作把过 $p$ 的一根轴转到另一根轴。轴共有 $n$ 根，每根轴可以转到另外 $n-1$ 根上去，但我们不应把“$y$ 转到 $x$”同“$x$ 转到 $y$”算作不同的转动，所以独立转动的总数是 $\frac{1}{2}n(n-1)$。于是我们有
\begin{equation*}
n + \frac{1}{2}n(n-1) = \frac{1}{2}n(n+1)
\tag{3.189}
\end{equation*}
个 $\mathbf{R}^n$ 的独立对称性。但我们的计数论证只涉及对称性在 $p$ 的邻域内的行为，而不是全局地遍及整个流形；所以即使存在曲率，这个计数也应相同。如果度规号差不是欧几里得型的，某些转动实际上会是推动（boost），但同样计数不变。等距变换的数目当然就是线性无关的 Killing 矢量场的数目。因此，我们把具有 $\frac{1}{2}n(n+1)$ 个 Killing 矢量的 $n$ 维流形称为\textbf{最大对称空间}（maximally symmetric space）。最大对称空间最熟悉的例子是 $n$ 维欧几里得空间 $\mathbf{R}^n$ 和 $n$ 维球面 $S^n$。对 $n$ 维球面，我们通常把等距变换看作由 $\frac{1}{2}n(n+1)$ 个独立转动组成，而不是转动和平移的某种混合集合。然而，如果考虑这些转动作用在某个固定点 $p$ 上，稍想片刻就会使我们确信：整个集合可以分解为绕该点的 $\frac{1}{2}n(n-1)$ 个转动（保持 $p$ 不动），以及另外 $n$ 个沿各个方向移动 $p$ 的转动，正如在 $\mathbf{R}^n$ 中那样。
% ---- 书页 140 / p155 ----
如果一个流形是最大对称的，那么曲率处处相同（正如平移类的等距变换所表达的那样），并且沿每个方向都相同（正如转动类的等距变换所表达的那样）。因此，只要知道了最大对称空间在一点的曲率，我们就知道了它所有的曲率。事实上，可能的最大对称空间只有寥寥数种；对它们进行分类的，是标量曲率 $R$（它将处处为常数）、维数 $n$、度规号差，或许还有与整体拓扑有关的一些分立信息（例如把 $n$ 维环面同 $\mathbf{R}^n$ 区分开来，以及把大小不同的环面彼此区分开）。由此可知，我们应当能够由 Ricci 标量 $R$ 重建这类空间的整个 Riemann 张量；让我们看看这是如何做到的。

基本想法很简单：既然几何在所有方向上看起来都一样，曲率张量在所有方向上也应该看起来一样。这可能意味着什么呢？先在某点 $p$ 选取局部惯性坐标，使得 $g_{\hat\mu\hat\nu} = \eta_{\hat\mu\hat\nu}$。当然，局部惯性坐标并不唯一；例如，我们可以在 $p$ 点作一个洛伦兹变换，度规分量仍将保持为 $\eta_{\hat\mu\hat\nu}$。（所谓“作一个洛伦兹变换”，我们实际上指的是改变 $T_p$ 中的基矢；在弯曲时空中，这只在单个点上有意义，而不是在一个区域上。）既然几何是最大对称的，我们希望 Riemann 张量也是如此；也就是说，$R_{\hat\rho\hat\sigma\hat\mu\hat\nu}$ 的分量在洛伦兹变换下也不应改变，因为时空中没有任何从优的方向。但在洛伦兹变换下分量不变的张量是独一无二的几种——度规、Kronecker $\delta$，以及 Levi-Civita 张量。
% ---- 书页 141 / p156 ----
这意味着，在这些坐标下、在这个点上，$R_{\hat\rho\hat\sigma\hat\mu\hat\nu}$ 的分量将正比于一个由这些不变张量构造出来的张量。尝试匹配 Riemann 张量的对称性就会发现，唯一的可能是
\begin{equation*}
R_{\hat\rho\hat\sigma\hat\mu\hat\nu} \propto g_{\hat\rho\hat\mu}g_{\hat\sigma\hat\nu} - g_{\hat\rho\hat\nu}g_{\hat\sigma\hat\mu}.
\tag{3.190}
\end{equation*}

但这是一个完全张量性的关系，所以它在任何坐标系中都必须成立。我们是在单个点 $p$ 处论证这个关系的，但在最大对称空间中所有的点生而平等，所以它在任何其他点上也必定成立。比例常数很容易固定：把两边缩并两次［左边变成 $R$，右边为 $n(n-1)$］。我们最终得到一个在任何最大对称空间中、在任何点、在任何坐标系里都成立的方程：
\begin{equation*}
R_{\rho\sigma\mu\nu} = \frac{R}{n(n-1)}(g_{\rho\mu}g_{\sigma\nu} - g_{\rho\nu}g_{\sigma\mu}).
\tag{3.191}
\end{equation*}

同样地，如果 Riemann 张量满足这个条件（其中 $R$ 在流形上为常数），那么度规就是最大对称的。在二维中，发现 $R$ 是常数就足以证明空间是最大对称的，因为曲率只有一个独立分量。在更高维时你就得多费些力气了。

于是，在局域上（暂且不管整体拓扑的问题），给定维数和号差的最大对称空间完全由 $R$ 确定。这类空间的基本分类，就是看 $R$ 为正、为零还是为负，因为 $R$ 的大小代表着空间尺寸的整体标度。对欧几里得号差，平直的最大对称空间是平面或适当的高维推广，而正曲率的则是球面。负曲率的最大对称欧几里得空间是双曲面（hyperboloid），记作 $H^n$。这些空间比较不为人熟悉，因为即使是二维双曲面也无法等距地嵌入 $\mathbf{R}^3$。让我们简要考察一下这个二维双曲面。

表示 $H^2$（它与 $\mathbf{R}^2$ 有相同的拓扑）的方式有好几种。一种简单的方式是把它们表示为\textbf{庞加莱半平面}（Poincaré half-plane），它是带有坐标 $\{x, y\}$ 的二维区域中 $y > 0$ 的部分，其度规为
\begin{equation*}
ds^2 = \frac{a^2}{y^2}(\mathrm{d}x^2 + \mathrm{d}y^2).
\tag{3.192}
\end{equation*}

庞加莱半平面的几何当然不同于 $\mathbf{R}^2$ 上半部分的几何，尽管使用了相似的坐标。例如，我们可以计算一条从 $y_1$ 竖直地（$x$ 为常数）延伸到 $y_2$ 的线段的长度：
% ---- 书页 142 / p157 ----
\begin{align*}
\Delta s &= \int_{y_1}^{y_2}\sqrt{g_{\mu\nu}\frac{\mathrm{d}x^\mu}{\mathrm{d}y}\frac{\mathrm{d}x^\nu}{\mathrm{d}y}}\,\mathrm{d}y\\
&= a\int_{y_1}^{y_2}\frac{\mathrm{d}y}{y}\\
&= a\ln\left(\frac{y_2}{y_1}\right).
\tag{3.193}
\end{align*}

这根本不是我们在欧几里得空间中所预期的结果 $\Delta s = y_2 - y_1$。特别地，注意对于趋近边界 $y = 0$ 的路径，路径长度变为无穷。换句话说，它其实根本算不上什么边界；就生活在双曲面上的任何居民而言，它无限遥远。

式 (3.192) 的非零 Christoffel 符号为
\begin{align*}
\Gamma^x_{xy} &= \Gamma^x_{yx} = -y^{-1}\\
\Gamma^y_{xx} &= y^{-1}\\
\Gamma^y_{yy} &= -y^{-1}.
\tag{3.194}
\end{align*}
由此出发容易证明，测地线满足
\begin{equation*}
(x - x_0)^2 + y^2 = l^2,
\tag{3.195}
\end{equation*}
其中 $x_0$ 和 $l$ 为常数。这种形式的曲线是圆心位于 $x$ 轴上的半圆，如图 3.8 所示。在 $x_0 \to \infty$、$l \to \infty$ 而 $l - x_0$ 固定的极限下，我们得到一条竖直的直线。仿照我们在第 3.7 节末尾对 $S^2$ 的讨论，我们算得 Riemann 张量的一个代表性分量为
\begin{equation*}
R^x{}_{yxy} = -y^{-2}.
\tag{3.196}
\end{equation*}

%FIG{3.8}: {带负曲率度规的上半平面。测地线是与 $x$ 轴垂直相交的半圆和直线。}
与二维球面一样，所有其他分量要么为零，要么通过对称性与该分量相联系。这不过是反映了这样一个事实：我们正处在二维
% 本块末句在原书书页143顶部继续（"mensions, ..."），下一块请用 %JOIN% 续接本句。
